%%%%%%%%%%%%%%%%%%%%%%%%%%%%%%%%%%%%%%%%%%%%%%%%%%%%%%%%%%%%
% 7.7 FUNCTIONS OF RANDOM VARIABLES
% The Composition of Advanced Mathematics
%%%%%%%%%%%%%%%%%%%%%%%%%%%%%%%%%%%%%%%%%%%%%%%%%%%%%%%%%%%%

\section{Functions of Random Variables}
\label{sec:functions-random-variables}

\prerequisite{
Random variables, discrete and continuous distributions, joint
distributions, cumulative distribution functions, expectation,
multivariable calculus, derivatives, inverse functions, double
integrals, and Jacobians.
}

\learningobjectives{
By the end of this section, the reader should be able to:
\begin{itemize}
    \item construct new random variables from existing random variables;
    \item determine the support of a transformed random variable;
    \item transform discrete probability mass functions;
    \item use the CDF method for continuous transformations;
    \item transform densities under one-to-one monotone mappings;
    \item handle many-to-one transformations;
    \item use location and scale transformations;
    \item find distributions of sums and differences;
    \item compute discrete and continuous convolutions;
    \item derive distributions of maxima and minima;
    \item transform pairs and random vectors using Jacobians;
    \item determine joint densities under invertible transformations;
    \item find marginal distributions after multivariable transformations;
    \item derive distributions of ratios and products;
    \item use polar-coordinate transformations for random variables;
    \item recognize transformations that generate standard distributions;
    \item understand order statistics at an introductory level;
    \item derive distributions of sample minima and maxima;
    \item distinguish transformations of values from transformations of
          densities;
    \item connect transformation methods with simulation, statistics,
          geometry, and probability transforms.
\end{itemize}
}

%%%%%%%%%%%%%%%%%%%%%%%%%%%%%%%%%%%%%%%%%%%%%%%%%%%%%%%%%%%%
\subsection{Why Transform Random Variables?}
\label{subsec:why-transform-random-variables}
%%%%%%%%%%%%%%%%%%%%%%%%%%%%%%%%%%%%%%%%%%%%%%%%%%%%%%%%%%%%

A probability model often begins with random variables whose
distributions are known.

One then defines new quantities such as

\[
Y=g(X)
\]

or

\[
Z=g(X,Y).
\]

Examples include:

\[
Y=X^2,
\]

\[
Z=X+Y,
\]

\[
R=\frac{X}{Y},
\]

\[
M=\max(X_1,\ldots,X_n),
\]

and

\[
\overline{X}
=
\frac1n
\sum_{i=1}^{n}X_i.
\]

The central problem is:

\[
\boxed{
\text{Given the distribution of the original variables, find the
distribution of the transformed variable.}
}
\]

%%%%%%%%%%%%%%%%%%%%%%%%%%%%%%%%%%%%%%%%%%%%%%%%%%%%%%%%%%%%
\subsection{Functions of One Random Variable}
\label{subsec:function-one-rv}
%%%%%%%%%%%%%%%%%%%%%%%%%%%%%%%%%%%%%%%%%%%%%%%%%%%%%%%%%%%%

Let

\[
X:\Omega\to\mathbb{R}
\]

be a random variable.

If

\[
g:\mathbb{R}\to\mathbb{R},
\]

then

\[
\boxed{
Y=g(X)
}
\]

is another random variable.

Explicitly,

\[
Y(\omega)
=
g(X(\omega)).
\]

Thus the transformation acts on the numerical value produced by \(X\).

%%%%%%%%%%%%%%%%%%%%%%%%%%%%%%%%%%%%%%%%%%%%%%%%%%%%%%%%%%%%
\subsection{Transforming the Support}
\label{subsec:function-transform-support}
%%%%%%%%%%%%%%%%%%%%%%%%%%%%%%%%%%%%%%%%%%%%%%%%%%%%%%%%%%%%

Before calculating probabilities, determine which values of \(Y\) are
possible.

If

\[
Y=g(X),
\]

then

\[
\boxed{
\operatorname{supp}(Y)
\subseteq
g(\operatorname{supp}(X)).
}
\]

In ordinary discrete examples this is often equality.

Determining the transformed support first prevents impossible values
from being assigned positive probability.

%%%%%%%%%%%%%%%%%%%%%%%%%%%%%%%%%%%%%%%%%%%%%%%%%%%%%%%%%%%%
\subsection{Discrete Transformations}
\label{subsec:discrete-transformations}
%%%%%%%%%%%%%%%%%%%%%%%%%%%%%%%%%%%%%%%%%%%%%%%%%%%%%%%%%%%%

Suppose \(X\) is discrete and

\[
Y=g(X).
\]

Then

\[
\boxed{
P(Y=y)
=
\sum_{\{x:g(x)=y\}}
P(X=x).
}
\]

Thus the probabilities of all \(x\)-values mapping to the same \(y\)
must be added.

%%%%%%%%%%%%%%%%%%%%%%%%%%%%%%%%%%%%%%%%%%%%%%%%%%%%%%%%%%%%
\subsection{Worked Example: One-to-One Discrete Transformation}
\label{subsec:worked-one-to-one-discrete}
%%%%%%%%%%%%%%%%%%%%%%%%%%%%%%%%%%%%%%%%%%%%%%%%%%%%%%%%%%%%

\begin{workedexamplebox}{Linear Transformation of a Die}

Let \(X\) be a fair die result and define

\[
Y=2X+1.
\]

The possible values are

\[
3,5,7,9,11,13.
\]

Because the mapping is one-to-one,

\[
P(Y=2x+1)
=
P(X=x)
=
\frac16.
\]

\begin{answerbox}
\[
\boxed{
P(Y=y)=\frac16
}
\]

for

\[
y\in\{3,5,7,9,11,13\}.
\]
\end{answerbox}

\end{workedexamplebox}

%%%%%%%%%%%%%%%%%%%%%%%%%%%%%%%%%%%%%%%%%%%%%%%%%%%%%%%%%%%%
\subsection{Many-to-One Discrete Transformations}
\label{subsec:many-to-one-discrete-transformations}
%%%%%%%%%%%%%%%%%%%%%%%%%%%%%%%%%%%%%%%%%%%%%%%%%%%%%%%%%%%%

If several values of \(X\) produce the same transformed value, their
probabilities must be combined.

For example, if

\[
Y=X^2,
\]

then both

\[
X=a
\]

and

\[
X=-a
\]

produce

\[
Y=a^2.
\]

Thus

\[
\boxed{
P(Y=a^2)
=
P(X=a)+P(X=-a)
}
\]

whenever these are the only preimages.

%%%%%%%%%%%%%%%%%%%%%%%%%%%%%%%%%%%%%%%%%%%%%%%%%%%%%%%%%%%%
\subsection{Worked Example: Squaring a Discrete Variable}
\label{subsec:worked-square-discrete}
%%%%%%%%%%%%%%%%%%%%%%%%%%%%%%%%%%%%%%%%%%%%%%%%%%%%%%%%%%%%

\begin{workedexamplebox}{A Symmetric Discrete Distribution}

Suppose

\[
P(X=-2)=0.2,
\]

\[
P(X=-1)=0.3,
\]

\[
P(X=1)=0.1,
\]

and

\[
P(X=2)=0.4.
\]

Let

\[
Y=X^2.
\]

Then

\[
P(Y=1)
=
P(X=-1)+P(X=1)
=
0.4.
\]

Also,

\[
P(Y=4)
=
P(X=-2)+P(X=2)
=
0.6.
\]

\begin{answerbox}
\[
\boxed{
P(Y=1)=0.4,
\qquad
P(Y=4)=0.6.
}
\]
\end{answerbox}

\end{workedexamplebox}

%%%%%%%%%%%%%%%%%%%%%%%%%%%%%%%%%%%%%%%%%%%%%%%%%%%%%%%%%%%%
\subsection{The CDF Transformation Method}
\label{subsec:cdf-transformation-method}
%%%%%%%%%%%%%%%%%%%%%%%%%%%%%%%%%%%%%%%%%%%%%%%%%%%%%%%%%%%%

The CDF method is one of the most general techniques for transforming a
single random variable.

If

\[
Y=g(X),
\]

begin with

\[
\boxed{
F_Y(y)
=
P(Y\leq y).
}
\]

Then rewrite the event

\[
\{g(X)\leq y\}
\]

as an event involving \(X\).

Finally use the known distribution of \(X\).

%%%%%%%%%%%%%%%%%%%%%%%%%%%%%%%%%%%%%%%%%%%%%%%%%%%%%%%%%%%%
\subsection{Increasing Transformations}
\label{subsec:increasing-transformations}
%%%%%%%%%%%%%%%%%%%%%%%%%%%%%%%%%%%%%%%%%%%%%%%%%%%%%%%%%%%%

Suppose \(g\) is strictly increasing.

Then

\[
g(X)\leq y
\]

is equivalent to

\[
X\leq g^{-1}(y).
\]

Therefore,

\[
\boxed{
F_Y(y)
=
F_X(g^{-1}(y)).
}
\]

This is the simplest transformation case.

%%%%%%%%%%%%%%%%%%%%%%%%%%%%%%%%%%%%%%%%%%%%%%%%%%%%%%%%%%%%
\subsection{Decreasing Transformations}
\label{subsec:decreasing-transformations}
%%%%%%%%%%%%%%%%%%%%%%%%%%%%%%%%%%%%%%%%%%%%%%%%%%%%%%%%%%%%

Suppose \(g\) is strictly decreasing.

Then

\[
g(X)\leq y
\]

is equivalent to

\[
X\geq g^{-1}(y).
\]

For a continuous \(X\),

\[
F_Y(y)
=
P(X\geq g^{-1}(y)).
\]

Thus,

\[
\boxed{
F_Y(y)
=
1-
F_X(g^{-1}(y))
}
\]

when endpoint masses are absent.

In the general case, left limits may be needed.

%%%%%%%%%%%%%%%%%%%%%%%%%%%%%%%%%%%%%%%%%%%%%%%%%%%%%%%%%%%%
\subsection{Density Transformation for Monotone Functions}
\label{subsec:density-monotone-transform}
%%%%%%%%%%%%%%%%%%%%%%%%%%%%%%%%%%%%%%%%%%%%%%%%%%%%%%%%%%%%

Suppose \(X\) has density \(f_X\), and

\[
Y=g(X)
\]

where \(g\) is one-to-one, differentiable, and monotone.

Let

\[
x=g^{-1}(y).
\]

Then

\[
\boxed{
f_Y(y)
=
f_X(g^{-1}(y))
\left|
\frac{d}{dy}
g^{-1}(y)
\right|.
}
\]

The absolute derivative adjusts for stretching or compression of
intervals.

%%%%%%%%%%%%%%%%%%%%%%%%%%%%%%%%%%%%%%%%%%%%%%%%%%%%%%%%%%%%
\subsection{Why the Derivative Factor Appears}
\label{subsec:why-derivative-factor}
%%%%%%%%%%%%%%%%%%%%%%%%%%%%%%%%%%%%%%%%%%%%%%%%%%%%%%%%%%%%

Suppose a small interval around \(x\) has width

\[
dx.
\]

Under

\[
y=g(x),
\]

its transformed width is approximately

\[
dy
=
g'(x)\,dx.
\]

Probability must be preserved:

\[
f_X(x)\,dx
\approx
f_Y(y)\,dy.
\]

Thus,

\[
f_Y(y)
\approx
f_X(x)
\left|
\frac{dx}{dy}
\right|.
\]

This leads to

\[
\boxed{
f_Y(y)
=
f_X(x)
\left|
\frac{dx}{dy}
\right|.
}
\]

%%%%%%%%%%%%%%%%%%%%%%%%%%%%%%%%%%%%%%%%%%%%%%%%%%%%%%%%%%%%
\subsection{Worked Example: Scaling a Uniform Variable}
\label{subsec:worked-scale-uniform-transform}
%%%%%%%%%%%%%%%%%%%%%%%%%%%%%%%%%%%%%%%%%%%%%%%%%%%%%%%%%%%%

\begin{workedexamplebox}{Linear Transformation}

Let

\[
X\sim\operatorname{Uniform}(0,1)
\]

and define

\[
Y=3X+2.
\]

Then

\[
x
=
\frac{y-2}{3},
\]

and

\[
\frac{dx}{dy}
=
\frac13.
\]

Also,

\[
2\leq y\leq5.
\]

Therefore,

\[
f_Y(y)
=
1\cdot\frac13.
\]

\begin{answerbox}
\[
\boxed{
Y\sim\operatorname{Uniform}(2,5).
}
\]
\end{answerbox}

\end{workedexamplebox}

%%%%%%%%%%%%%%%%%%%%%%%%%%%%%%%%%%%%%%%%%%%%%%%%%%%%%%%%%%%%
\subsection{Location-Scale Transformation}
\label{subsec:location-scale-transformation}
%%%%%%%%%%%%%%%%%%%%%%%%%%%%%%%%%%%%%%%%%%%%%%%%%%%%%%%%%%%%

Suppose

\[
Y=aX+b,
\qquad
a\neq0.
\]

Then

\[
X
=
\frac{Y-b}{a}.
\]

The transformed density is

\[
\boxed{
f_Y(y)
=
\frac1{|a|}
f_X
\left(
\frac{y-b}{a}
\right).
}
\]

This is the general location-scale formula.

%%%%%%%%%%%%%%%%%%%%%%%%%%%%%%%%%%%%%%%%%%%%%%%%%%%%%%%%%%%%
\subsection{Normal Location-Scale Transformation}
\label{subsec:normal-location-scale-transform}
%%%%%%%%%%%%%%%%%%%%%%%%%%%%%%%%%%%%%%%%%%%%%%%%%%%%%%%%%%%%

If

\[
Z\sim N(0,1)
\]

and

\[
X=\mu+\sigma Z,
\qquad
\sigma>0,
\]

then

\[
\boxed{
X\sim N(\mu,\sigma^2).
}
\]

Conversely,

\[
\boxed{
Z
=
\frac{X-\mu}{\sigma}
\sim
N(0,1).
}
\]

Thus standardization is a location-scale transformation.

%%%%%%%%%%%%%%%%%%%%%%%%%%%%%%%%%%%%%%%%%%%%%%%%%%%%%%%%%%%%
\subsection{Worked Example: Transforming an Exponential Variable}
\label{subsec:worked-scaled-exponential}
%%%%%%%%%%%%%%%%%%%%%%%%%%%%%%%%%%%%%%%%%%%%%%%%%%%%%%%%%%%%

\begin{workedexamplebox}{Scaling an Exponential Variable}

Suppose

\[
X\sim\operatorname{Exp}(\lambda)
\]

and define

\[
Y=cX,
\qquad
c>0.
\]

Then

\[
x=\frac{y}{c},
\]

and

\[
\frac{dx}{dy}
=
\frac1c.
\]

Therefore,

\[
f_Y(y)
=
\lambda
e^{-\lambda y/c}
\frac1c.
\]

Hence,

\begin{answerbox}
\[
\boxed{
Y
\sim
\operatorname{Exp}
\left(
\frac{\lambda}{c}
\right).
}
\]
\end{answerbox}

\end{workedexamplebox}

%%%%%%%%%%%%%%%%%%%%%%%%%%%%%%%%%%%%%%%%%%%%%%%%%%%%%%%%%%%%
\subsection{Many-to-One Continuous Transformations}
\label{subsec:many-to-one-continuous}
%%%%%%%%%%%%%%%%%%%%%%%%%%%%%%%%%%%%%%%%%%%%%%%%%%%%%%%%%%%%

Suppose

\[
Y=g(X)
\]

and several \(x\)-values satisfy

\[
g(x)=y.
\]

If the inverse branches are

\[
x_1(y),\ldots,x_m(y),
\]

then under suitable smoothness conditions,

\[
\boxed{
f_Y(y)
=
\sum_{i=1}^{m}
f_X(x_i(y))
\left|
\frac{dx_i}{dy}
\right|.
}
\]

Each inverse branch contributes density.

%%%%%%%%%%%%%%%%%%%%%%%%%%%%%%%%%%%%%%%%%%%%%%%%%%%%%%%%%%%%
\subsection{Worked Example: Squaring a Symmetric Variable}
\label{subsec:worked-square-continuous}
%%%%%%%%%%%%%%%%%%%%%%%%%%%%%%%%%%%%%%%%%%%%%%%%%%%%%%%%%%%%

\begin{workedexamplebox}{Transformation \(Y=X^2\)}

Suppose \(X\) has density \(f_X\) on the real line and define

\[
Y=X^2.
\]

For \(y>0\),

\[
x=\sqrt{y}
\]

or

\[
x=-\sqrt{y}.
\]

Also,

\[
\left|
\frac{dx}{dy}
\right|
=
\frac1{2\sqrt{y}}.
\]

Therefore,

\[
\begin{answerbox}
\boxed{
f_Y(y)
=
\frac{
f_X(\sqrt{y})
+
f_X(-\sqrt{y})
}{
2\sqrt{y}
},
\qquad
y>0.
}
\]
\end{answerbox}

\end{workedexamplebox}

%%%%%%%%%%%%%%%%%%%%%%%%%%%%%%%%%%%%%%%%%%%%%%%%%%%%%%%%%%%%
\subsection{Square of a Standard Normal Variable}
\label{subsec:square-standard-normal}
%%%%%%%%%%%%%%%%%%%%%%%%%%%%%%%%%%%%%%%%%%%%%%%%%%%%%%%%%%%%

Let

\[
Z\sim N(0,1)
\]

and define

\[
Y=Z^2.
\]

Then the transformation above gives a density of gamma type.

The resulting distribution is

\[
\boxed{
Z^2
\sim
\chi^2_1.
}
\]

This is the first connection between normal and chi-square variables.

%%%%%%%%%%%%%%%%%%%%%%%%%%%%%%%%%%%%%%%%%%%%%%%%%%%%%%%%%%%%
\subsection{Functions of Two Random Variables}
\label{subsec:functions-two-random-variables-full}
%%%%%%%%%%%%%%%%%%%%%%%%%%%%%%%%%%%%%%%%%%%%%%%%%%%%%%%%%%%%

Suppose

\[
X
\]

and

\[
Y
\]

have a known joint distribution.

Define

\[
\boxed{
Z=g(X,Y).
}
\]

Examples include

\[
Z=X+Y,
\]

\[
Z=X-Y,
\]

\[
Z=XY,
\]

\[
Z=\frac{X}{Y},
\]

and

\[
Z=\max(X,Y).
\]

Several techniques are available depending on the form of \(g\).

%%%%%%%%%%%%%%%%%%%%%%%%%%%%%%%%%%%%%%%%%%%%%%%%%%%%%%%%%%%%
\subsection{Distribution of a Discrete Sum}
\label{subsec:discrete-sum-convolution}
%%%%%%%%%%%%%%%%%%%%%%%%%%%%%%%%%%%%%%%%%%%%%%%%%%%%%%%%%%%%

Let

\[
Z=X+Y.
\]

For discrete random variables,

\[
\boxed{
P(Z=z)
=
\sum_x
P(X=x,Y=z-x).
}
\]

If \(X\) and \(Y\) are independent,

\[
\boxed{
P(Z=z)
=
\sum_x
p_X(x)p_Y(z-x).
}
\]

This operation is called convolution.

%%%%%%%%%%%%%%%%%%%%%%%%%%%%%%%%%%%%%%%%%%%%%%%%%%%%%%%%%%%%
\subsection{Discrete Convolution}
\label{subsec:discrete-convolution}
%%%%%%%%%%%%%%%%%%%%%%%%%%%%%%%%%%%%%%%%%%%%%%%%%%%%%%%%%%%%

For PMFs \(p_X\) and \(p_Y\), define

\[
\boxed{
(p_X*p_Y)(z)
=
\sum_x
p_X(x)p_Y(z-x).
}
\]

Then for independent \(X\) and \(Y\),

\[
\boxed{
p_{X+Y}
=
p_X*p_Y.
}
\]

The symbol

\[
*
\]

denotes convolution.

%%%%%%%%%%%%%%%%%%%%%%%%%%%%%%%%%%%%%%%%%%%%%%%%%%%%%%%%%%%%
\subsection{Worked Example: Sum of Two Bernoulli Variables}
\label{subsec:worked-sum-bernoulli}
%%%%%%%%%%%%%%%%%%%%%%%%%%%%%%%%%%%%%%%%%%%%%%%%%%%%%%%%%%%%

\begin{workedexamplebox}{Two Independent Bernoulli Variables}

Let

\[
X,Y
\sim
\operatorname{Bernoulli}(p)
\]

independently.

Then

\[
Z=X+Y.
\]

Possible values are

\[
0,1,2.
\]

We obtain

\[
P(Z=0)
=
(1-p)^2,
\]

\[
P(Z=1)
=
2p(1-p),
\]

and

\[
P(Z=2)
=
p^2.
\]

\begin{answerbox}
\[
\boxed{
Z
\sim
\operatorname{Bin}(2,p).
}
\]
\end{answerbox}

\end{workedexamplebox}

%%%%%%%%%%%%%%%%%%%%%%%%%%%%%%%%%%%%%%%%%%%%%%%%%%%%%%%%%%%%
\subsection{Sum of Independent Binomial Variables}
\label{subsec:sum-binomial-variables}
%%%%%%%%%%%%%%%%%%%%%%%%%%%%%%%%%%%%%%%%%%%%%%%%%%%%%%%%%%%%

Suppose

\[
X\sim\operatorname{Bin}(m,p)
\]

and

\[
Y\sim\operatorname{Bin}(n,p)
\]

independently.

Then

\[
X+Y
\]

counts successes across

\[
m+n
\]

independent Bernoulli trials with common success probability \(p\).

Therefore,

\[
\boxed{
X+Y
\sim
\operatorname{Bin}(m+n,p).
}
\]

%%%%%%%%%%%%%%%%%%%%%%%%%%%%%%%%%%%%%%%%%%%%%%%%%%%%%%%%%%%%
\subsection{Sum of Independent Poisson Variables}
\label{subsec:sum-poisson-functions}
%%%%%%%%%%%%%%%%%%%%%%%%%%%%%%%%%%%%%%%%%%%%%%%%%%%%%%%%%%%%

If

\[
X\sim\operatorname{Poisson}(\lambda_1)
\]

and

\[
Y\sim\operatorname{Poisson}(\lambda_2)
\]

independently, then

\[
\boxed{
X+Y
\sim
\operatorname{Poisson}(\lambda_1+\lambda_2).
}
\]

This can be proved by convolution or by probability generating
functions.

%%%%%%%%%%%%%%%%%%%%%%%%%%%%%%%%%%%%%%%%%%%%%%%%%%%%%%%%%%%%
\subsection{Continuous Convolution}
\label{subsec:continuous-convolution}
%%%%%%%%%%%%%%%%%%%%%%%%%%%%%%%%%%%%%%%%%%%%%%%%%%%%%%%%%%%%

Suppose \(X\) and \(Y\) are independent continuous random variables.

For

\[
Z=X+Y,
\]

the density is

\[
\boxed{
f_Z(z)
=
\int_{-\infty}^{\infty}
f_X(x)
f_Y(z-x)
\,dx.
}
\]

Equivalently,

\[
\boxed{
f_Z
=
f_X*f_Y.
}
\]

%%%%%%%%%%%%%%%%%%%%%%%%%%%%%%%%%%%%%%%%%%%%%%%%%%%%%%%%%%%%
\subsection{Deriving Continuous Convolution}
\label{subsec:derive-continuous-convolution}
%%%%%%%%%%%%%%%%%%%%%%%%%%%%%%%%%%%%%%%%%%%%%%%%%%%%%%%%%%%%

Let

\[
Z=X+Y.
\]

Consider the transformation

\[
U=X,
\qquad
Z=X+Y.
\]

Then

\[
Y=Z-U.
\]

The joint density becomes

\[
f_{U,Z}(u,z)
=
f_{X,Y}(u,z-u).
\]

For independent variables,

\[
=
f_X(u)f_Y(z-u).
\]

Integrating over \(u\),

\[
\boxed{
f_Z(z)
=
\int_{-\infty}^{\infty}
f_X(u)f_Y(z-u)\,du.
}
\]

%%%%%%%%%%%%%%%%%%%%%%%%%%%%%%%%%%%%%%%%%%%%%%%%%%%%%%%%%%%%
\subsection{Worked Example: Sum of Two Uniform Variables}
\label{subsec:worked-sum-uniform}
%%%%%%%%%%%%%%%%%%%%%%%%%%%%%%%%%%%%%%%%%%%%%%%%%%%%%%%%%%%%

\begin{workedexamplebox}{Triangular Density from Convolution}

Let

\[
X,Y
\sim
\operatorname{Uniform}(0,1)
\]

independently.

Let

\[
Z=X+Y.
\]

For

\[
0<z<1,
\]

the valid range is

\[
0<x<z.
\]

Therefore,

\[
f_Z(z)
=
\int_0^z
1\,dx
=
z.
\]

For

\[
1\leq z<2,
\]

the valid range is

\[
z-1<x<1.
\]

Thus,

\[
f_Z(z)
=
\int_{z-1}^{1}
1\,dx
=
2-z.
\]

\begin{answerbox}
\[
\boxed{
f_Z(z)
=
\begin{cases}
z,&0<z<1,\\
2-z,&1\leq z<2,\\
0,&\text{otherwise}.
\end{cases}
}
\]
\end{answerbox}

\end{workedexamplebox}

%%%%%%%%%%%%%%%%%%%%%%%%%%%%%%%%%%%%%%%%%%%%%%%%%%%%%%%%%%%%
\subsection{Sums of Independent Normal Variables}
\label{subsec:sum-independent-normal}
%%%%%%%%%%%%%%%%%%%%%%%%%%%%%%%%%%%%%%%%%%%%%%%%%%%%%%%%%%%%

Suppose

\[
X\sim N(\mu_1,\sigma_1^2)
\]

and

\[
Y\sim N(\mu_2,\sigma_2^2)
\]

independently.

Then

\[
\boxed{
X+Y
\sim
N
\left(
\mu_1+\mu_2,
\,
\sigma_1^2+\sigma_2^2
\right).
}
\]

More generally, linear combinations of jointly normal random variables
remain normal.

%%%%%%%%%%%%%%%%%%%%%%%%%%%%%%%%%%%%%%%%%%%%%%%%%%%%%%%%%%%%
\subsection{Difference of Independent Variables}
\label{subsec:difference-independent-variables}
%%%%%%%%%%%%%%%%%%%%%%%%%%%%%%%%%%%%%%%%%%%%%%%%%%%%%%%%%%%%

If

\[
Z=X-Y,
\]

then write

\[
Z=X+(-Y).
\]

Thus the distribution may be obtained by transforming \(Y\) to \(-Y\)
and then convolving.

For independent normal variables,

\[
\boxed{
X-Y
\sim
N
\left(
\mu_X-\mu_Y,
\,
\sigma_X^2+\sigma_Y^2
\right).
}
\]

%%%%%%%%%%%%%%%%%%%%%%%%%%%%%%%%%%%%%%%%%%%%%%%%%%%%%%%%%%%%
\subsection{Sample Sums}
\label{subsec:sample-sums}
%%%%%%%%%%%%%%%%%%%%%%%%%%%%%%%%%%%%%%%%%%%%%%%%%%%%%%%%%%%%

For random variables

\[
X_1,\ldots,X_n,
\]

define

\[
\boxed{
S_n
=
X_1+\cdots+X_n.
}
\]

If the variables are independent, the distribution of \(S_n\) is the
iterated convolution of their distributions.

If they are i.i.d.,

\[
\boxed{
f_{S_n}
=
f_X^{*n}
}
\]

in the continuous case.

%%%%%%%%%%%%%%%%%%%%%%%%%%%%%%%%%%%%%%%%%%%%%%%%%%%%%%%%%%%%
\subsection{Sample Mean as a Transformation}
\label{subsec:sample-mean-transformation}
%%%%%%%%%%%%%%%%%%%%%%%%%%%%%%%%%%%%%%%%%%%%%%%%%%%%%%%%%%%%

The sample mean is

\[
\boxed{
\overline{X}
=
\frac1n
\sum_{i=1}^{n}X_i.
}
\]

Thus its distribution can be obtained from the distribution of

\[
S_n
=
\sum_{i=1}^{n}X_i
\]

through the scale transformation

\[
\overline{X}
=
\frac{S_n}{n}.
\]

This becomes fundamental in statistical inference.

%%%%%%%%%%%%%%%%%%%%%%%%%%%%%%%%%%%%%%%%%%%%%%%%%%%%%%%%%%%%
\subsection{Normal Sample Means}
\label{subsec:normal-sample-means}
%%%%%%%%%%%%%%%%%%%%%%%%%%%%%%%%%%%%%%%%%%%%%%%%%%%%%%%%%%%%

If

\[
X_1,\ldots,X_n
\]

are independent and

\[
X_i\sim N(\mu,\sigma^2),
\]

then

\[
S_n
\sim
N(n\mu,n\sigma^2).
\]

Therefore,

\[
\boxed{
\overline{X}
\sim
N
\left(
\mu,
\frac{\sigma^2}{n}
\right).
}
\]

This identity is exact for normal samples.

%%%%%%%%%%%%%%%%%%%%%%%%%%%%%%%%%%%%%%%%%%%%%%%%%%%%%%%%%%%%
\subsection{Maximum of Random Variables}
\label{subsec:maximum-random-variables}
%%%%%%%%%%%%%%%%%%%%%%%%%%%%%%%%%%%%%%%%%%%%%%%%%%%%%%%%%%%%

Let

\[
M=\max(X_1,\ldots,X_n).
\]

The event

\[
M\leq m
\]

occurs exactly when

\[
X_1\leq m,
\ldots,
X_n\leq m.
\]

Therefore,

\[
\boxed{
F_M(m)
=
P(X_1\leq m,\ldots,X_n\leq m).
}
\]

%%%%%%%%%%%%%%%%%%%%%%%%%%%%%%%%%%%%%%%%%%%%%%%%%%%%%%%%%%%%
\subsection{Maximum of Independent Variables}
\label{subsec:maximum-independent-rvs}
%%%%%%%%%%%%%%%%%%%%%%%%%%%%%%%%%%%%%%%%%%%%%%%%%%%%%%%%%%%%

If

\[
X_1,\ldots,X_n
\]

are independent, then

\[
\boxed{
F_M(m)
=
\prod_{i=1}^{n}
F_{X_i}(m).
}
\]

If they are i.i.d. with common CDF \(F\),

\[
\boxed{
F_M(m)
=
[F(m)]^n.
}
\]

%%%%%%%%%%%%%%%%%%%%%%%%%%%%%%%%%%%%%%%%%%%%%%%%%%%%%%%%%%%%
\subsection{Density of an i.i.d. Maximum}
\label{subsec:density-iid-maximum}
%%%%%%%%%%%%%%%%%%%%%%%%%%%%%%%%%%%%%%%%%%%%%%%%%%%%%%%%%%%%

Suppose the variables are i.i.d. continuous with common density \(f\)
and CDF \(F\).

Differentiating,

\[
\boxed{
f_M(m)
=
n
[F(m)]^{n-1}
f(m).
}
\]

This is the density of the sample maximum.

%%%%%%%%%%%%%%%%%%%%%%%%%%%%%%%%%%%%%%%%%%%%%%%%%%%%%%%%%%%%
\subsection{Minimum of Random Variables}
\label{subsec:minimum-random-variables}
%%%%%%%%%%%%%%%%%%%%%%%%%%%%%%%%%%%%%%%%%%%%%%%%%%%%%%%%%%%%

Let

\[
L=\min(X_1,\ldots,X_n).
\]

The event

\[
L>l
\]

occurs exactly when every variable exceeds \(l\).

Therefore, for independent variables,

\[
P(L>l)
=
\prod_{i=1}^{n}
P(X_i>l).
\]

Hence,

\[
\boxed{
F_L(l)
=
1-
\prod_{i=1}^{n}
[1-F_{X_i}(l)].
}
\]

%%%%%%%%%%%%%%%%%%%%%%%%%%%%%%%%%%%%%%%%%%%%%%%%%%%%%%%%%%%%
\subsection{Minimum of i.i.d. Variables}
\label{subsec:minimum-iid}
%%%%%%%%%%%%%%%%%%%%%%%%%%%%%%%%%%%%%%%%%%%%%%%%%%%%%%%%%%%%

If the variables are i.i.d. with CDF \(F\),

\[
\boxed{
F_L(l)
=
1-[1-F(l)]^n.
}
\]

If \(F\) has density \(f\), then

\[
\boxed{
f_L(l)
=
n
[1-F(l)]^{n-1}
f(l).
}
\]

%%%%%%%%%%%%%%%%%%%%%%%%%%%%%%%%%%%%%%%%%%%%%%%%%%%%%%%%%%%%
\subsection{Worked Example: Maximum of Uniform Variables}
\label{subsec:worked-maximum-uniform}
%%%%%%%%%%%%%%%%%%%%%%%%%%%%%%%%%%%%%%%%%%%%%%%%%%%%%%%%%%%%

\begin{workedexamplebox}{Maximum of \(n\) Uniform Variables}

Let

\[
X_1,\ldots,X_n
\]

be i.i.d.

\[
\operatorname{Uniform}(0,1).
\]

For

\[
0\leq m\leq1,
\]

\[
F(m)=m.
\]

Therefore,

\[
F_M(m)=m^n.
\]

Differentiating,

\[
\begin{answerbox}
\[
\boxed{
f_M(m)
=
nm^{n-1},
\qquad
0<m<1.
}
\]
\end{answerbox}

\end{workedexamplebox}

%%%%%%%%%%%%%%%%%%%%%%%%%%%%%%%%%%%%%%%%%%%%%%%%%%%%%%%%%%%%
\subsection{Order Statistics}
\label{subsec:order-statistics}
%%%%%%%%%%%%%%%%%%%%%%%%%%%%%%%%%%%%%%%%%%%%%%%%%%%%%%%%%%%%

Suppose

\[
X_1,\ldots,X_n
\]

are random variables.

After arranging the observed values from smallest to largest, write

\[
\boxed{
X_{(1)}
\leq
X_{(2)}
\leq
\cdots
\leq
X_{(n)}.
}
\]

These are called order statistics.

Here,

\[
X_{(1)}
=
\min_i X_i
\]

and

\[
X_{(n)}
=
\max_i X_i.
\]

%%%%%%%%%%%%%%%%%%%%%%%%%%%%%%%%%%%%%%%%%%%%%%%%%%%%%%%%%%%%
\subsection{Density of an Order Statistic}
\label{subsec:density-order-statistic}
%%%%%%%%%%%%%%%%%%%%%%%%%%%%%%%%%%%%%%%%%%%%%%%%%%%%%%%%%%%%

Suppose

\[
X_1,\ldots,X_n
\]

are i.i.d. continuous with common CDF \(F\) and density \(f\).

Then the density of the \(k\)-th order statistic is

\[
\boxed{
f_{X_{(k)}}(x)
=
\frac{
n!
}{
(k-1)!(n-k)!
}
[F(x)]^{k-1}
[1-F(x)]^{n-k}
f(x).
}
\]

This formula counts which observations fall below and above \(x\).

%%%%%%%%%%%%%%%%%%%%%%%%%%%%%%%%%%%%%%%%%%%%%%%%%%%%%%%%%%%%
\subsection{Uniform Order Statistics}
\label{subsec:uniform-order-statistics}
%%%%%%%%%%%%%%%%%%%%%%%%%%%%%%%%%%%%%%%%%%%%%%%%%%%%%%%%%%%%

For i.i.d.

\[
\operatorname{Uniform}(0,1)
\]

variables,

\[
F(x)=x
\]

and

\[
f(x)=1.
\]

Therefore,

\[
\boxed{
f_{X_{(k)}}(x)
=
\frac{
n!
}{
(k-1)!(n-k)!
}
x^{k-1}
(1-x)^{n-k},
}
\]

for

\[
0<x<1.
\]

Thus,

\[
\boxed{
X_{(k)}
\sim
\operatorname{Beta}
(k,n-k+1).
}
\]

%%%%%%%%%%%%%%%%%%%%%%%%%%%%%%%%%%%%%%%%%%%%%%%%%%%%%%%%%%%%
\subsection{Transformation of Two Continuous Variables}
\label{subsec:two-variable-transformations}
%%%%%%%%%%%%%%%%%%%%%%%%%%%%%%%%%%%%%%%%%%%%%%%%%%%%%%%%%%%%

Suppose

\[
(U,V)
=
T(X,Y)
\]

where

\[
U=g_1(X,Y),
\qquad
V=g_2(X,Y).
\]

If the transformation is one-to-one, differentiable, and invertible,
write the inverse as

\[
X=x(u,v),
\qquad
Y=y(u,v).
\]

Then the joint density transforms using a Jacobian.

%%%%%%%%%%%%%%%%%%%%%%%%%%%%%%%%%%%%%%%%%%%%%%%%%%%%%%%%%%%%
\subsection{Jacobian Determinant}
\label{subsec:jacobian-determinant-probability}
%%%%%%%%%%%%%%%%%%%%%%%%%%%%%%%%%%%%%%%%%%%%%%%%%%%%%%%%%%%%

The Jacobian matrix of the inverse transformation is

\[
\frac{\partial(x,y)}{\partial(u,v)}
=
\begin{pmatrix}
\dfrac{\partial x}{\partial u}
&
\dfrac{\partial x}{\partial v}
\\[3mm]
\dfrac{\partial y}{\partial u}
&
\dfrac{\partial y}{\partial v}
\end{pmatrix}.
\]

Its determinant is

\[
\boxed{
J
=
\det
\left(
\frac{\partial(x,y)}{\partial(u,v)}
\right).
}
\]

The density transformation uses

\[
|J|.
\]

%%%%%%%%%%%%%%%%%%%%%%%%%%%%%%%%%%%%%%%%%%%%%%%%%%%%%%%%%%%%
\subsection{Joint Change-of-Variables Formula}
\label{subsec:joint-change-of-variables}
%%%%%%%%%%%%%%%%%%%%%%%%%%%%%%%%%%%%%%%%%%%%%%%%%%%%%%%%%%%%

\begin{theorem}[Two-Dimensional Change of Variables]
Suppose the transformation between \((X,Y)\) and \((U,V)\) is
one-to-one and differentiable with a differentiable inverse.

Then

\[
\boxed{
f_{U,V}(u,v)
=
f_{X,Y}
\bigl(
x(u,v),
y(u,v)
\bigr)
\left|
\det
\frac{\partial(x,y)}{\partial(u,v)}
\right|.
}
\]
\end{theorem}

The transformed support must also be determined.

%%%%%%%%%%%%%%%%%%%%%%%%%%%%%%%%%%%%%%%%%%%%%%%%%%%%%%%%%%%%
\subsection{Worked Example: Sum and Original Variable}
\label{subsec:worked-jacobian-sum}
%%%%%%%%%%%%%%%%%%%%%%%%%%%%%%%%%%%%%%%%%%%%%%%%%%%%%%%%%%%%

\begin{workedexamplebox}{Transformation \(U=X+Y,\;V=X\)}

Define

\[
U=X+Y,
\qquad
V=X.
\]

Then

\[
X=V,
\]

and

\[
Y=U-V.
\]

The inverse Jacobian is

\[
\frac{\partial(x,y)}{\partial(u,v)}
=
\begin{pmatrix}
0&1\\
1&-1
\end{pmatrix}.
\]

Its determinant is

\[
-1.
\]

Therefore,

\[
|J|=1.
\]

Hence,

\begin{answerbox}
\[
\boxed{
f_{U,V}(u,v)
=
f_{X,Y}(v,u-v).
}
\]
\end{answerbox}

\end{workedexamplebox}

%%%%%%%%%%%%%%%%%%%%%%%%%%%%%%%%%%%%%%%%%%%%%%%%%%%%%%%%%%%%
\subsection{Recovering a Single Transformed Density}
\label{subsec:marginalizing-transformed-density}
%%%%%%%%%%%%%%%%%%%%%%%%%%%%%%%%%%%%%%%%%%%%%%%%%%%%%%%%%%%%

After finding

\[
f_{U,V}(u,v),
\]

the marginal density of \(U\) is

\[
\boxed{
f_U(u)
=
\int
f_{U,V}(u,v)\,dv.
}
\]

This technique is often easier than trying to find \(f_U\) directly.

It is particularly useful for:

\[
\text{sums},
\quad
\text{ratios},
\quad
\text{products},
\quad
\text{quadratic forms}.
\]

%%%%%%%%%%%%%%%%%%%%%%%%%%%%%%%%%%%%%%%%%%%%%%%%%%%%%%%%%%%%
\subsection{Distribution of a Ratio}
\label{subsec:distribution-ratio}
%%%%%%%%%%%%%%%%%%%%%%%%%%%%%%%%%%%%%%%%%%%%%%%%%%%%%%%%%%%%

Suppose

\[
R=\frac{X}{Y}
\]

with \(Y\neq0\).

A useful transformation is

\[
R=\frac{X}{Y},
\qquad
V=Y.
\]

Then

\[
X=RV,
\qquad
Y=V.
\]

The inverse Jacobian determinant has absolute value

\[
\boxed{
|V|.
}
\]

Therefore,

\[
\boxed{
f_{R,V}(r,v)
=
f_{X,Y}(rv,v)|v|.
}
\]

Integrating over \(v\),

\[
\boxed{
f_R(r)
=
\int_{-\infty}^{\infty}
f_{X,Y}(rv,v)|v|\,dv.
}
\]

%%%%%%%%%%%%%%%%%%%%%%%%%%%%%%%%%%%%%%%%%%%%%%%%%%%%%%%%%%%%
\subsection{Ratio of Independent Standard Normals}
\label{subsec:ratio-normal-cauchy}
%%%%%%%%%%%%%%%%%%%%%%%%%%%%%%%%%%%%%%%%%%%%%%%%%%%%%%%%%%%%

Let

\[
X,Y
\]

be independent standard normal random variables.

Then

\[
\boxed{
\frac{X}{Y}
}
\]

has a standard Cauchy distribution.

Thus,

\[
\boxed{
\frac{X}{Y}
\sim
\operatorname{Cauchy}(0,1).
}
\]

This is a major structural relationship between the normal and Cauchy
families.

%%%%%%%%%%%%%%%%%%%%%%%%%%%%%%%%%%%%%%%%%%%%%%%%%%%%%%%%%%%%
\subsection{Distribution of a Product}
\label{subsec:distribution-product}
%%%%%%%%%%%%%%%%%%%%%%%%%%%%%%%%%%%%%%%%%%%%%%%%%%%%%%%%%%%%

Suppose

\[
U=XY.
\]

A useful transformation is

\[
U=XY,
\qquad
V=X.
\]

Then

\[
X=V,
\]

\[
Y=\frac{U}{V}.
\]

The inverse Jacobian has absolute determinant

\[
\boxed{
\frac1{|V|}.
}
\]

Therefore,

\[
\boxed{
f_U(u)
=
\int_{-\infty}^{\infty}
f_{X,Y}
\left(
v,\frac{u}{v}
\right)
\frac1{|v|}
\,dv,
}
\]

over the valid region.

%%%%%%%%%%%%%%%%%%%%%%%%%%%%%%%%%%%%%%%%%%%%%%%%%%%%%%%%%%%%
\subsection{Polar Transformations}
\label{subsec:polar-transformations-probability}
%%%%%%%%%%%%%%%%%%%%%%%%%%%%%%%%%%%%%%%%%%%%%%%%%%%%%%%%%%%%

For two variables, an important transformation is

\[
X=R\cos\Theta,
\]

\[
Y=R\sin\Theta.
\]

Then

\[
R
=
\sqrt{X^2+Y^2},
\]

and

\[
\Theta
=
\operatorname{atan2}(Y,X).
\]

The Jacobian determinant is

\[
\boxed{
\left|
\frac{\partial(x,y)}{\partial(r,\theta)}
\right|
=
r.
}
\]

%%%%%%%%%%%%%%%%%%%%%%%%%%%%%%%%%%%%%%%%%%%%%%%%%%%%%%%%%%%%
\subsection{Worked Example: Independent Standard Normals in Polar Form}
\label{subsec:worked-polar-standard-normals}
%%%%%%%%%%%%%%%%%%%%%%%%%%%%%%%%%%%%%%%%%%%%%%%%%%%%%%%%%%%%

\begin{workedexamplebox}{Radius of a Bivariate Standard Normal}

Let

\[
X,Y
\]

be independent standard normal random variables.

Their joint density is

\[
f_{X,Y}(x,y)
=
\frac1{2\pi}
e^{-(x^2+y^2)/2}.
\]

Use

\[
x=r\cos\theta,
\qquad
y=r\sin\theta.
\]

Since

\[
x^2+y^2=r^2
\]

and the Jacobian is \(r\),

\[
f_{R,\Theta}(r,\theta)
=
\frac1{2\pi}
e^{-r^2/2}r.
\]

For

\[
r\geq0,
\qquad
0\leq\theta<2\pi.
\]

Integrating over \(\theta\),

\[
f_R(r)
=
re^{-r^2/2}.
\]

\begin{answerbox}
\[
\boxed{
f_R(r)
=
re^{-r^2/2},
\qquad
r\geq0.
}
\]
\end{answerbox}

\end{workedexamplebox}

%%%%%%%%%%%%%%%%%%%%%%%%%%%%%%%%%%%%%%%%%%%%%%%%%%%%%%%%%%%%
\subsection{Rayleigh Distribution}
\label{subsec:rayleigh-distribution}
%%%%%%%%%%%%%%%%%%%%%%%%%%%%%%%%%%%%%%%%%%%%%%%%%%%%%%%%%%%%

The radial variable in the previous example has a Rayleigh
distribution.

A Rayleigh density with scale \(\sigma>0\) is

\[
\boxed{
f_R(r)
=
\frac{r}{\sigma^2}
\exp
\left(
-\frac{r^2}{2\sigma^2}
\right),
\qquad
r\geq0.
}
\]

When

\[
\sigma=1,
\]

this becomes

\[
re^{-r^2/2}.
\]

%%%%%%%%%%%%%%%%%%%%%%%%%%%%%%%%%%%%%%%%%%%%%%%%%%%%%%%%%%%%
\subsection{Chi-Square from Normal Squares}
\label{subsec:chi-square-normal-squares}
%%%%%%%%%%%%%%%%%%%%%%%%%%%%%%%%%%%%%%%%%%%%%%%%%%%%%%%%%%%%

Let

\[
Z_1,\ldots,Z_\nu
\]

be independent standard normal variables.

Then

\[
\boxed{
Q
=
Z_1^2+\cdots+Z_\nu^2
\sim
\chi^2_\nu.
}
\]

Thus the chi-square distribution arises naturally from the squared
Euclidean length of a standard normal random vector.

%%%%%%%%%%%%%%%%%%%%%%%%%%%%%%%%%%%%%%%%%%%%%%%%%%%%%%%%%%%%
\subsection{Student's \(t\)-Distribution Construction}
\label{subsec:t-distribution-construction}
%%%%%%%%%%%%%%%%%%%%%%%%%%%%%%%%%%%%%%%%%%%%%%%%%%%%%%%%%%%%

Let

\[
Z\sim N(0,1)
\]

and

\[
V\sim\chi^2_\nu
\]

independently.

Then

\[
\boxed{
T
=
\frac{
Z
}{
\sqrt{V/\nu}
}
}
\]

has Student's \(t\)-distribution with \(\nu\) degrees of freedom.

We write

\[
\boxed{
T\sim t_\nu.
}
\]

This construction becomes fundamental in statistical inference.

%%%%%%%%%%%%%%%%%%%%%%%%%%%%%%%%%%%%%%%%%%%%%%%%%%%%%%%%%%%%
\subsection{\(F\)-Distribution Construction}
\label{subsec:f-distribution-construction}
%%%%%%%%%%%%%%%%%%%%%%%%%%%%%%%%%%%%%%%%%%%%%%%%%%%%%%%%%%%%

Suppose

\[
U\sim\chi^2_{\nu_1}
\]

and

\[
V\sim\chi^2_{\nu_2}
\]

independently.

Then

\[
\boxed{
F
=
\frac{
U/\nu_1
}{
V/\nu_2
}
}
\]

has an \(F\)-distribution with

\[
\nu_1
\]

and

\[
\nu_2
\]

degrees of freedom.

We write

\[
\boxed{
F\sim F_{\nu_1,\nu_2}.
}
\]

%%%%%%%%%%%%%%%%%%%%%%%%%%%%%%%%%%%%%%%%%%%%%%%%%%%%%%%%%%%%
\subsection{Transformation of Random Vectors}
\label{subsec:random-vector-transformations}
%%%%%%%%%%%%%%%%%%%%%%%%%%%%%%%%%%%%%%%%%%%%%%%%%%%%%%%%%%%%

For an \(n\)-dimensional random vector

\[
\mathbf{X},
\]

consider

\[
\mathbf{Y}
=
g(\mathbf{X}).
\]

If \(g\) is invertible and differentiable, then

\[
\boxed{
f_{\mathbf{Y}}(\mathbf{y})
=
f_{\mathbf{X}}
(g^{-1}(\mathbf{y}))
\left|
\det
Dg^{-1}(\mathbf{y})
\right|.
}
\]

Here

\[
Dg^{-1}
\]

is the Jacobian matrix of the inverse transformation.

%%%%%%%%%%%%%%%%%%%%%%%%%%%%%%%%%%%%%%%%%%%%%%%%%%%%%%%%%%%%
\subsection{Linear Transformations of Random Vectors}
\label{subsec:linear-transform-random-vectors}
%%%%%%%%%%%%%%%%%%%%%%%%%%%%%%%%%%%%%%%%%%%%%%%%%%%%%%%%%%%%

Suppose

\[
\mathbf{Y}
=
A\mathbf{X}+\mathbf{b}.
\]

Then expectation transforms as

\[
\boxed{
\mathbb{E}[\mathbf{Y}]
=
A\mathbb{E}[\mathbf{X}]
+
\mathbf{b}.
}
\]

If \(\mathbf{X}\) has covariance matrix \(\Sigma_X\), then

\[
\boxed{
\Sigma_Y
=
A\Sigma_XA^T.
}
\]

This formula is central in multivariate probability and statistics.

%%%%%%%%%%%%%%%%%%%%%%%%%%%%%%%%%%%%%%%%%%%%%%%%%%%%%%%%%%%%
\subsection{Linear Transformation of a Multivariate Normal}
\label{subsec:linear-transform-multivariate-normal}
%%%%%%%%%%%%%%%%%%%%%%%%%%%%%%%%%%%%%%%%%%%%%%%%%%%%%%%%%%%%

If

\[
\mathbf{X}
\sim
N_n(\boldsymbol{\mu},\Sigma)
\]

and

\[
\mathbf{Y}
=
A\mathbf{X}+\mathbf{b},
\]

then

\[
\boxed{
\mathbf{Y}
\sim
N_m
\left(
A\boldsymbol{\mu}+\mathbf{b},
A\Sigma A^T
\right),
}
\]

provided the dimensions are compatible.

Thus multivariate normal distributions are closed under linear
transformations.

%%%%%%%%%%%%%%%%%%%%%%%%%%%%%%%%%%%%%%%%%%%%%%%%%%%%%%%%%%%%
\subsection{Probability Integral Transform}
\label{subsec:probability-integral-transform}
%%%%%%%%%%%%%%%%%%%%%%%%%%%%%%%%%%%%%%%%%%%%%%%%%%%%%%%%%%%%

Suppose \(X\) has a continuous CDF \(F_X\).

Then

\[
\boxed{
U=F_X(X)
}
\]

has a uniform distribution on

\[
(0,1).
\]

Thus,

\[
\boxed{
U\sim\operatorname{Uniform}(0,1).
}
\]

This is called the probability integral transform.

%%%%%%%%%%%%%%%%%%%%%%%%%%%%%%%%%%%%%%%%%%%%%%%%%%%%%%%%%%%%
\subsection{Inverse Probability Transform}
\label{subsec:inverse-probability-transform}
%%%%%%%%%%%%%%%%%%%%%%%%%%%%%%%%%%%%%%%%%%%%%%%%%%%%%%%%%%%%

Conversely, let

\[
U\sim\operatorname{Uniform}(0,1).
\]

If \(F\) is a continuous strictly increasing CDF, then

\[
\boxed{
X=F^{-1}(U)
}
\]

has CDF \(F\).

This provides the theoretical basis for inverse-transform simulation.

%%%%%%%%%%%%%%%%%%%%%%%%%%%%%%%%%%%%%%%%%%%%%%%%%%%%%%%%%%%%
\subsection{Worked Example: Generating an Exponential Variable}
\label{subsec:worked-generate-exponential}
%%%%%%%%%%%%%%%%%%%%%%%%%%%%%%%%%%%%%%%%%%%%%%%%%%%%%%%%%%%%

\begin{workedexamplebox}{Inverse Transform Sampling}

Let

\[
U\sim\operatorname{Uniform}(0,1).
\]

For an exponential distribution,

\[
F(x)
=
1-e^{-\lambda x}.
\]

Set

\[
U
=
1-e^{-\lambda X}.
\]

Solving for \(X\),

\[
X
=
-\frac1\lambda
\ln(1-U).
\]

Since

\[
1-U
\sim
\operatorname{Uniform}(0,1),
\]

one may write

\begin{answerbox}
\[
\boxed{
X
=
-\frac1\lambda
\ln U.
}
\]
\end{answerbox}

\end{workedexamplebox}

%%%%%%%%%%%%%%%%%%%%%%%%%%%%%%%%%%%%%%%%%%%%%%%%%%%%%%%%%%%%
\subsection{Transformations and Expectations}
\label{subsec:transformations-expectation}
%%%%%%%%%%%%%%%%%%%%%%%%%%%%%%%%%%%%%%%%%%%%%%%%%%%%%%%%%%%%

Finding the full distribution of

\[
Y=g(X)
\]

is unnecessary if one only wants

\[
\mathbb{E}[Y].
\]

By the law of the unconscious statistician,

\[
\boxed{
\mathbb{E}[g(X)]
=
\int
g(x)f_X(x)\,dx
}
\]

for continuous \(X\), or

\[
\boxed{
\mathbb{E}[g(X)]
=
\sum_x
g(x)p_X(x)
}
\]

for discrete \(X\).

Thus distribution transformation and expectation transformation are
different tasks.

%%%%%%%%%%%%%%%%%%%%%%%%%%%%%%%%%%%%%%%%%%%%%%%%%%%%%%%%%%%%
\subsection{Transformation Strategy}
\label{subsec:transformation-strategy}
%%%%%%%%%%%%%%%%%%%%%%%%%%%%%%%%%%%%%%%%%%%%%%%%%%%%%%%%%%%%

When faced with a transformed random variable, ask:

\begin{enumerate}
    \item Is the original variable discrete or continuous?
    \item What is the transformed support?
    \item Is the transformation one-to-one?
    \item Is it monotone?
    \item Would the CDF method be simpler?
    \item Are there multiple inverse branches?
    \item Is the new variable a sum?
    \item Can convolution be used?
    \item Is a multivariable transformation needed?
    \item What Jacobian appears?
    \item Must one marginalize after transformation?
\end{enumerate}

%%%%%%%%%%%%%%%%%%%%%%%%%%%%%%%%%%%%%%%%%%%%%%%%%%%%%%%%%%%%
\subsection{Common Errors}
\label{subsec:functions-random-variables-common-errors}
%%%%%%%%%%%%%%%%%%%%%%%%%%%%%%%%%%%%%%%%%%%%%%%%%%%%%%%%%%%%

\subsubsection{Forgetting the Transformed Support}

A correct density formula applied outside its valid support is still an
incorrect probability model.

Always transform the support.

%%%%%%%%%%%%%%%%%%%%%%%%%%%%%%%%%%%%%%%%%%%%%%%%%%%%%%%%%%%%
\subsubsection{Forgetting Multiple Preimages}

For a many-to-one transformation, all inverse branches must be included.

For example,

\[
Y=X^2
\]

typically has both

\[
X=\sqrt{Y}
\]

and

\[
X=-\sqrt{Y}.
\]

%%%%%%%%%%%%%%%%%%%%%%%%%%%%%%%%%%%%%%%%%%%%%%%%%%%%%%%%%%%%
\subsubsection{Forgetting the Absolute Value of the Derivative}

The one-dimensional transformation formula uses

\[
\boxed{
\left|
\frac{dx}{dy}
\right|.
}
\]

Without the absolute value, a decreasing transformation could produce a
negative density.

%%%%%%%%%%%%%%%%%%%%%%%%%%%%%%%%%%%%%%%%%%%%%%%%%%%%%%%%%%%%
\subsubsection{Using the Forward Derivative Instead of the Inverse Derivative}

If the formula is written as

\[
f_Y(y)
=
f_X(g^{-1}(y))
\left|
\frac{d}{dy}g^{-1}(y)
\right|,
\]

the derivative belongs to the inverse transformation.

%%%%%%%%%%%%%%%%%%%%%%%%%%%%%%%%%%%%%%%%%%%%%%%%%%%%%%%%%%%%
\subsubsection{Using Convolution Without Independence}

The formula

\[
f_{X+Y}(z)
=
\int
f_X(x)f_Y(z-x)\,dx
\]

assumes independence.

Without independence, the joint density must be used.

%%%%%%%%%%%%%%%%%%%%%%%%%%%%%%%%%%%%%%%%%%%%%%%%%%%%%%%%%%%%
\subsubsection{Ignoring Convolution Limits}

Even if the integral is formally written over the whole real line, the
actual nonzero range is determined by the supports of \(X\) and \(Y\).

%%%%%%%%%%%%%%%%%%%%%%%%%%%%%%%%%%%%%%%%%%%%%%%%%%%%%%%%%%%%
\subsubsection{Forgetting the Jacobian}

A multivariable transformation changes area or volume.

The Jacobian determinant corrects the density for this distortion.

%%%%%%%%%%%%%%%%%%%%%%%%%%%%%%%%%%%%%%%%%%%%%%%%%%%%%%%%%%%%
\subsubsection{Using the Wrong Jacobian Direction}

If using

\[
f_{U,V}(u,v)
=
f_{X,Y}(x(u,v),y(u,v))
\left|
\frac{\partial(x,y)}{\partial(u,v)}
\right|,
\]

the Jacobian is for the inverse mapping.

%%%%%%%%%%%%%%%%%%%%%%%%%%%%%%%%%%%%%%%%%%%%%%%%%%%%%%%%%%%%
\subsubsection{Ignoring Non-One-to-One Transformations}

If a multivariable transformation is not globally one-to-one, the
domain may need to be partitioned into regions on which the mapping is
one-to-one.

%%%%%%%%%%%%%%%%%%%%%%%%%%%%%%%%%%%%%%%%%%%%%%%%%%%%%%%%%%%%
\subsubsection{Confusing Distribution of a Sum with Sum of Densities}

In general,

\[
f_{X+Y}(z)
\neq
f_X(z)+f_Y(z).
\]

Independent sums are obtained using convolution.

%%%%%%%%%%%%%%%%%%%%%%%%%%%%%%%%%%%%%%%%%%%%%%%%%%%%%%%%%%%%
\subsubsection{Assuming Means Determine Transformed Distributions}

Knowing

\[
\mathbb{E}[X]
\]

and

\[
\operatorname{Var}(X)
\]

is generally not enough to determine the distribution of

\[
g(X).
\]

The full distribution may be required.

%%%%%%%%%%%%%%%%%%%%%%%%%%%%%%%%%%%%%%%%%%%%%%%%%%%%%%%%%%%%
\subsection{Applications}
\label{subsec:functions-random-variables-applications}
%%%%%%%%%%%%%%%%%%%%%%%%%%%%%%%%%%%%%%%%%%%%%%%%%%%%%%%%%%%%

\begin{applicationbox}

\textbf{Statistics.}
Sample means, sample variances, \(t\)-statistics, \(F\)-statistics, and
order statistics are functions of random samples.

\medskip

\textbf{Simulation.}
Inverse-transform methods construct nonuniform random variables from
uniform random numbers.

\medskip

\textbf{Reliability.}
System lifetime may be the minimum or maximum of component lifetimes.

\medskip

\textbf{Engineering.}
Derived quantities such as power, energy, stress, and signal magnitude
are functions of uncertain inputs.

\medskip

\textbf{Finance.}
Portfolio return is a linear combination of random asset returns.

\medskip

\textbf{Physics.}
Coordinate transformations and radial variables naturally require
Jacobian methods.

\medskip

\textbf{Queueing Theory.}
Waiting and service times are often sums of random variables.

\medskip

\textbf{Extreme-Value Analysis.}
Sample maxima and minima are fundamental transformed variables.

\medskip

\textbf{Multivariate Statistics.}
Linear transformations of random vectors determine covariance
structures and principal directions.

\end{applicationbox}

%%%%%%%%%%%%%%%%%%%%%%%%%%%%%%%%%%%%%%%%%%%%%%%%%%%%%%%%%%%%
\subsection{Exercises}
\label{subsec:functions-random-variables-exercises}
%%%%%%%%%%%%%%%%%%%%%%%%%%%%%%%%%%%%%%%%%%%%%%%%%%%%%%%%%%%%

\begin{exercise}[1]
Let \(X\) take values

\[
-2,-1,1,2
\]

with equal probabilities.

Find the PMF of

\[
Y=X^2.
\]
\end{exercise}

\begin{exercise}[1]
Let

\[
Y=3X+4.
\]

If

\[
\operatorname{supp}(X)=\{0,1,2,3\},
\]

find the support of \(Y\).
\end{exercise}

\begin{exercise}[1]
State the one-dimensional density transformation formula for a
monotone transformation.
\end{exercise}

\begin{exercise}[1]
Define convolution.
\end{exercise}

\begin{exercise}[1]
State the convolution formula for two independent discrete random
variables.
\end{exercise}

\begin{exercise}[1]
State the convolution formula for two independent continuous random
variables.
\end{exercise}

\begin{exercise}[1]
State the CDF of the maximum of \(n\) i.i.d. variables.
\end{exercise}

\begin{exercise}[1]
State the CDF of the minimum of \(n\) i.i.d. variables.
\end{exercise}

\begin{exercise}[1]
Define an order statistic.
\end{exercise}

\begin{exercise}[1]
State the two-dimensional Jacobian transformation formula.
\end{exercise}

\begin{exercise}[2]
Let

\[
X\sim\operatorname{Uniform}(0,2)
\]

and define

\[
Y=5X.
\]

Find the distribution of \(Y\).
\end{exercise}

\begin{exercise}[2]
Let

\[
X\sim\operatorname{Exp}(2)
\]

and define

\[
Y=4X.
\]

Find the distribution of \(Y\).
\end{exercise}

\begin{exercise}[2]
Suppose \(X\) has density

\[
f_X(x)
=
\frac12,
\qquad
-1<x<1.
\]

Find the density of

\[
Y=X^2.
\]
\end{exercise}

\begin{exercise}[2]
Let \(X\) and \(Y\) be independent Bernoulli\((p)\) variables.

Find the PMF of

\[
X+Y.
\]
\end{exercise}

\begin{exercise}[2]
Let

\[
X\sim\operatorname{Poisson}(2)
\]

and

\[
Y\sim\operatorname{Poisson}(5)
\]

independently.

Find the distribution of \(X+Y\).
\end{exercise}

\begin{exercise}[2]
Let

\[
X\sim N(3,4)
\]

and

\[
Y\sim N(5,9)
\]

independently.

Find the distribution of

\[
X+Y.
\]
\end{exercise}

\begin{exercise}[2]
For the previous variables, find the distribution of

\[
X-Y.
\]
\end{exercise}

\begin{exercise}[2]
Let \(X,Y\) be independent

\[
\operatorname{Uniform}(0,1).
\]

Find the density of

\[
X+Y.
\]
\end{exercise}

\begin{exercise}[2]
Let \(X_1,X_2,X_3\) be i.i.d. uniform on \([0,1]\).

Find the CDF of

\[
M=\max(X_1,X_2,X_3).
\]
\end{exercise}

\begin{exercise}[2]
Find the density of \(M\) in the previous exercise.
\end{exercise}

\begin{exercise}[2]
For the same sample, find the CDF of

\[
L=\min(X_1,X_2,X_3).
\]
\end{exercise}

\begin{exercise}[3]
Derive the one-dimensional density transformation formula using the
CDF method.
\end{exercise}

\begin{exercise}[3]
Derive the many-to-one density formula for

\[
Y=X^2.
\]
\end{exercise}

\begin{exercise}[3]
If

\[
X\sim N(0,1),
\]

show that

\[
X^2\sim\chi^2_1.
\]
\end{exercise}

\begin{exercise}[3]
Derive the discrete convolution formula for \(X+Y\).
\end{exercise}

\begin{exercise}[3]
Derive the continuous convolution formula using a Jacobian
transformation.
\end{exercise}

\begin{exercise}[3]
Use convolution to prove that the sum of independent gamma variables
with common rate and shapes \(\alpha_1,\alpha_2\) is gamma with shape

\[
\alpha_1+\alpha_2.
\]
\end{exercise}

\begin{exercise}[3]
Use convolution to prove that independent normal variables are closed
under addition.
\end{exercise}

\begin{exercise}[3]
Let

\[
X_1,\ldots,X_n
\]

be i.i.d. exponential variables with rate \(\lambda\).

Find the distribution of their minimum.
\end{exercise}

\begin{exercise}[3]
Show that the minimum in the previous exercise is exponential with rate

\[
n\lambda.
\]
\end{exercise}

\begin{exercise}[3]
Derive the density of the \(k\)-th order statistic for an i.i.d.
continuous sample.
\end{exercise}

\begin{exercise}[3]
Show that if

\[
X_1,\ldots,X_n
\sim
\operatorname{Uniform}(0,1)
\]

i.i.d., then

\[
X_{(k)}
\sim
\operatorname{Beta}(k,n-k+1).
\]
\end{exercise}

\begin{exercise}[3]
For

\[
U=X+Y,
\qquad
V=X-Y,
\]

find the inverse transformation and the absolute Jacobian determinant.
\end{exercise}

\begin{exercise}[3]
Suppose \(X,Y\) are independent standard normal variables.

Transform to

\[
R=\sqrt{X^2+Y^2},
\qquad
\Theta=\operatorname{atan2}(Y,X),
\]

and derive the joint density of \((R,\Theta)\).
\end{exercise}

\begin{exercise}[3]
Use the previous result to derive the Rayleigh distribution of \(R\).
\end{exercise}

\begin{exercise}[3]
Let \(X,Y\) be independent standard normal variables.

Show that

\[
X/Y
\]

has a standard Cauchy distribution.
\end{exercise}

\begin{exercise}[3]
If

\[
\mathbf{Y}
=
A\mathbf{X}+\mathbf{b},
\]

prove

\[
\operatorname{Cov}(\mathbf{Y})
=
A\operatorname{Cov}(\mathbf{X})A^T.
\]
\end{exercise}

\begin{exercise}[3]
Prove the probability integral transform for a continuous strictly
increasing CDF.
\end{exercise}

\begin{exercise}[3]
Use inverse-transform sampling to construct a

\[
\operatorname{Uniform}(a,b)
\]

variable from a

\[
\operatorname{Uniform}(0,1)
\]

variable.
\end{exercise}

\begin{exercise}[4]
Derive the density of Student's \(t\)-distribution from its normal and
chi-square representation.
\end{exercise}

\begin{exercise}[4]
Derive the density of the \(F\)-distribution from independent
chi-square variables.
\end{exercise}

\begin{exercise}[4]
Study transformations that produce beta distributions from gamma
variables.
\end{exercise}

\begin{exercise}[4]
Show that if

\[
X\sim\operatorname{Gamma}(\alpha,\lambda)
\]

and

\[
Y\sim\operatorname{Gamma}(\beta,\lambda)
\]

independently, then

\[
\frac{X}{X+Y}
\]

has a beta distribution.
\end{exercise}

\begin{exercise}[4]
Investigate the independence of

\[
\frac{X}{X+Y}
\]

and

\[
X+Y
\]

for independent gamma variables with common rate.
\end{exercise}

\begin{exercise}[4]
Study the Box--Muller transformation for generating independent standard
normal random variables from independent uniform variables.
\end{exercise}

\begin{exercise}[4]
Investigate transformation methods for noninvertible mappings in higher
dimensions.
\end{exercise}

\begin{exercise}[4]
Study order statistics and derive the joint density of two order
statistics.
\end{exercise}

\begin{exercise}[4]
Investigate sample ranges

\[
R=X_{(n)}-X_{(1)}
\]

and derive their distribution for simple parent distributions.
\end{exercise}

\begin{exercise}[4]
Study extreme-value limits for transformed sample maxima.
\end{exercise}

%%%%%%%%%%%%%%%%%%%%%%%%%%%%%%%%%%%%%%%%%%%%%%%%%%%%%%%%%%%%
\subsection{Conceptual Summary}
\label{subsec:functions-random-variables-summary}
%%%%%%%%%%%%%%%%%%%%%%%%%%%%%%%%%%%%%%%%%%%%%%%%%%%%%%%%%%%%

If

\[
Y=g(X),
\]

then the transformed support should be determined first.

For a discrete variable,

\[
\boxed{
P(Y=y)
=
\sum_{\{x:g(x)=y\}}
P(X=x).
}
\]

For a monotone continuous transformation,

\[
\boxed{
f_Y(y)
=
f_X(g^{-1}(y))
\left|
\frac{d}{dy}
g^{-1}(y)
\right|.
}
\]

For multiple inverse branches,

\[
\boxed{
f_Y(y)
=
\sum_i
f_X(x_i(y))
\left|
\frac{dx_i}{dy}
\right|.
}
\]

For independent discrete variables,

\[
\boxed{
p_{X+Y}(z)
=
\sum_x
p_X(x)p_Y(z-x).
}
\]

For independent continuous variables,

\[
\boxed{
f_{X+Y}(z)
=
\int_{-\infty}^{\infty}
f_X(x)f_Y(z-x)\,dx.
}
\]

For an i.i.d. sample maximum,

\[
\boxed{
F_{X_{(n)}}(x)
=
[F(x)]^n.
}
\]

For an i.i.d. sample minimum,

\[
\boxed{
F_{X_{(1)}}(x)
=
1-[1-F(x)]^n.
}
\]

The \(k\)-th order statistic has density

\[
\boxed{
f_{X_{(k)}}(x)
=
\frac{
n!
}{
(k-1)!(n-k)!
}
[F(x)]^{k-1}
[1-F(x)]^{n-k}
f(x).
}
\]

For a two-dimensional invertible transformation,

\[
\boxed{
f_{U,V}(u,v)
=
f_{X,Y}
(x(u,v),y(u,v))
\left|
\det
\frac{\partial(x,y)}{\partial(u,v)}
\right|.
}
\]

For linear random-vector transformations,

\[
\boxed{
\mathbb{E}[A\mathbf{X}+\mathbf{b}]
=
A\mathbb{E}[\mathbf{X}]+\mathbf{b},
}
\]

and

\[
\boxed{
\operatorname{Cov}(A\mathbf{X}+\mathbf{b})
=
A\Sigma A^T.
}
\]

Functions of random variables therefore connect

\[
\boxed{
\text{probability}
+
\text{calculus}
+
\text{geometry}
+
\text{convolution}
+
\text{simulation}.
}
\]

%%%%%%%%%%%%%%%%%%%%%%%%%%%%%%%%%%%%%%%%%%%%%%%%%%%%%%%%%%%%
\subsection{Section Connections}
\label{subsec:functions-random-variables-connections}
%%%%%%%%%%%%%%%%%%%%%%%%%%%%%%%%%%%%%%%%%%%%%%%%%%%%%%%%%%%%

\chapterconnection{
Functions of Random Variables builds on the distributional framework of
Sections~\ref{sec:random-variables} through
\ref{sec:expectation-moments}. It provides the transformation,
convolution, Jacobian, order-statistic, and random-vector tools required
to derive new probability distributions from known ones. The next
section, Probability Transforms, provides an algebraic alternative to
direct convolution by encoding distributions with generating,
moment-generating, Laplace, and characteristic functions. Later, the
Laws of Large Numbers and Central Limit Theorems study transformations
of large sums and averages, while Statistics repeatedly uses sample
means, sample variances, order statistics, chi-square variables,
Student's \(t\)-statistics, and \(F\)-statistics derived through the
methods developed here.
}

%%%%%%%%%%%%%%%%%%%%%%%%%%%%%%%%%%%%%%%%%%%%%%%%%%%%%%%%%%%%
% SECTION SOURCE TRACKING
%%%%%%%%%%%%%%%%%%%%%%%%%%%%%%%%%%%%%%%%%%%%%%%%%%%%%%%%%%%%

%------------------------------------------------------------
% 7.7 Functions of Random Variables
%------------------------------------------------------------
%
% Source basis:
%   - Standard probability transformation theory.
%   - Standard convolution methods.
%   - Standard multivariable change-of-variables theory.
%   - Standard order-statistic theory.
%
% Used for:
%   - functions of random variables;
%   - transformed supports;
%   - discrete transformations;
%   - one-to-one transformations;
%   - many-to-one transformations;
%   - CDF transformation method;
%   - density transformation formula;
%   - location-scale transformations;
%   - normal standardization;
%   - exponential scaling;
%   - discrete convolution;
%   - continuous convolution;
%   - sums of Bernoulli variables;
%   - sums of binomial variables;
%   - sums of Poisson variables;
%   - sums and differences of normal variables;
%   - sample sums;
%   - sample means;
%   - maxima and minima;
%   - order statistics;
%   - uniform order statistics;
%   - two-variable transformations;
%   - Jacobian determinants;
%   - joint change of variables;
%   - ratios;
%   - products;
%   - polar transformations;
%   - Rayleigh distribution;
%   - chi-square construction;
%   - Student's t construction;
%   - F-distribution construction;
%   - random-vector transformations;
%   - multivariate normal transformations;
%   - probability integral transform;
%   - inverse-transform sampling;
%   - applications and exercises.
%
% Notes:
%   - Probability transforms are developed next in Section 7.8.
%   - Student's t and F distributions are developed more fully in
%     Statistics.
%   - Order statistics reappear in estimation and nonparametric
%     statistics.
%   - Jacobians connect directly with multivariable calculus and
%     multivariate probability.
%   - Extreme-value limits are developed only as advanced previews here.
%
%%%%%%%%%%%%%%%%%%%%%%%%%%%%%%%%%%%%%%%%%%%%%%%%%%%%%%%%%%%%